\begingroup
\subsection{Sensitivity to Measurement Noise and External Forces}
\hl{Fig.~\ref{fig:robustness} summarizes the two frozen policies and their arithmetic mean. Panel (d) displays the combined position-noise sweep over $\sigma_p=0.75$--$1.00$ mm; this range was restricted after evaluation, with the two higher evaluated levels retained in the experiment archive. Without perturbations, mean success was $96.65\%$ on this test set. With position noise, success was $97.60\%$ at $\sigma_p=0.5$ mm, but fell to $53.10\%$ at $1$ mm and $0.05\%$ at $2$ mm. Thus, noise on the scale of the position-success tolerance substantially degraded reaching performance. Under the tested nonzero orientation-noise levels ($0.0025$--$0.020$ rad per component), mean success ranged from $97.30\%$ to $99.15\%$. Small nonmonotonic increases do not establish a general benefit from measurement noise.}

\hl{Force-only mean success remained $96.65\%$ at 0, 5, and 10 N and $96.60\%$ at 20 N. This characterizes the policy together with its underlying velocity servo; it does not isolate disturbance rejection by the learned policy. With simultaneous 20-N force and position noise, mean success decreased from $96.25\%$ at $0.75$ mm to $88.60\%$, $80.60\%$, $67.05\%$, and $51.70\%$ at $0.85$, $0.90$, $0.95$, and $1.00$ mm, respectively. At $1.00$ mm, the individual policies achieved $35.8\%$ and $67.6\%$, exposing sensitivity to the training seed.}
\par\endgroup
